\startSection
\selectlanguage{french}
\sectionTitle{L'Hymne National Français}

\begin{question}{Qui a composé \enquote{La Marseillaise}, devenue l'hymne national de la France?}
  \choice[correct]{Claude Joseph Rouget de Lisle}
  \choice{Jean-Baptiste Lully}
  \choice{Hector Berlioz}
  \choice{François-Joseph Gossec}
  \choice{Jean-Philippe Rameau}
\end{question}

\begin{question}{Dans quelle ville le chant qui deviendra \enquote{La Marseillaise} a-t-il été composé en 1792?}
  \choice{Paris}
  \choice{Marseille}
  \choice{Lyon}
  \choice[correct]{Strasbourg}
  \choice{Bordeaux}
\end{question}

\begin{question}{Quel était le titre original du chant avant qu'il ne soit connu sous le nom de \enquote{La Marseillaise}?}
  \choice{Chant de guerre de l'armée française}
  \choice{Hymne des Marseillais}
  \choice[correct]{Chant de guerre pour l'armée du Rhin}
  \choice{Chant patriotique de la République}
  \choice{Marche des volontaires}
\end{question}

\begin{question}{En quelle année \enquote{La Marseillaise} a-t-elle été adoptée comme hymne national pendant la Révolution française?}
  \choice{1792}
  \choice[correct]{1795}
  \choice{1804}
  \choice{1815}
  \choice{1848}
\end{question}

\begin{question}{Pourquoi ce chant porte-t-il le nom de \enquote{La Marseillaise}, bien qu'il ait été composé à Strasbourg?}
  \choice{Il célèbre une victoire militaire à Marseille.}
  \choice{Le compositeur était originaire de Marseille.}
  \choice[correct]{Des fédérés marseillais l'ont popularisé à Paris.}
  \choice{Il fut imprimé pour la première fois à Marseille.}
  \choice{Il remplaça un hymne municipal marseillais.}
\end{question}

\begin{question}{Dans son contexte révolutionnaire, le refrain de \enquote{La Marseillaise} sert surtout à produire quel effet civique?}
  \choice{Présenter une prière liturgique.}
  \choice{Décrire une cérémonie de cour.}
  \choice{Énumérer les provinces françaises.}
  \choice[correct]{Appeler les citoyens à la mobilisation collective.}
  \choice{Expliquer les règles d'un scrutin national.}
\end{question}
\shuffleSection
\renderSection